\chapter{Error reporting}

\section{The chain}
\estab\ None of the three documented errors is raised as the number 40, 41 or 42. The
SysEx engine keeps a parse record whose field~4 is a status byte. At the end of a
bulk-dump session that status is mapped through a table to a message id, which selects a
display list.

\begin{center}
\begin{tabular}{lll}
\toprule
message id & screen & raised by status \\
\midrule
\bytes{0x23} & \textsc{completed!} & \bytes{0x00} \\
\bytes{0x22} & \textsc{error 40!} & \bytes{0x08}, \bytes{0x09}, \bytes{0x0A} \\
\bytes{0x21} & \textsc{error 41!} & everything not listed elsewhere \\
\bytes{0x20} & \textsc{error 42!} & \bytes{0x20} (and \bytes{0x18}, \bytes{0x1F}, never raised) \\
\bytes{0x0F} & \textsc{error 21!} ``Memory full'' & \bytes{0x16} \\
\bottomrule
\end{tabular}
\end{center}

\begin{tcolorbox}[colback=red!4!white,colframe=red!50!black,title=\textbf{Not every
bulk-dump failure is a 4\textit{x} error}]
\estab\ Status \bytes{0x16} --- the incoming dump is larger than the destination area ---
produces \textsc{error 21! Memory full}. A document that assumes every SysEx bulk-dump
failure is 40, 41 or 42 gets this one wrong.
\end{tcolorbox}

\section{ERROR 40 --- wrong identification}
\estab\ Raised when the grammar fails at depth 2, 3 or 4, that is at message byte 3, 4
or 5. For a message addressed to the WSA1 (\bytes{0x2D}) this means byte 3 was not
\bytes{0x04}, or byte 4 was not \bytes{0x00}/\bytes{0x01}, or byte 5 was not \bytes{0x11}.

\section{ERROR 41 --- reception}
\estab\ The catch-all. The conditions actually raised include: first byte of a message in
a session is not \bytes{0xF0}; second byte is not \bytes{0x50}; a status byte other than
\bytes{0xF7} arrived inside an open message; the UART reported a framing, parity or
overrun error; a reply timed out; the message exceeded 254 bytes; an odd body length; the
continuation flag was neither \bytes{0x00} nor \bytes{0x01}; the checksum did not match.

\section{ERROR 42 --- transmission}
\estab\ \textbf{Exactly one reachable cause.} Status \bytes{0x20} is raised in the loop
that streams a dump out, when the inter-processor link reports that a block transfer is
still outstanding about 1.02\,s after the wait began. Statuses \bytes{0x18} and
\bytes{0x1F} also map to ERROR 42 but are written nowhere in any of the four ROM images.

\section{Where errors are visible}
\estab\ Status bytes are written whenever a complete message is parsed, from either MIDI
ring, but the popup appears only from the \textsc{sysex bulk dump} screen. A foreign SysEx
arriving on any other screen sets no visible error.

\estab\ While a bulk-dump session is open, unrelated SysEx traffic on the MIDI input can
set the status and so colour the popup the session ends with.

\section{Verbatim screen text}
\estab\ The firmware's own spelling is reproduced here as it appears:

\begin{quote}\small
\textsc{error 40!} --- ``The Identification (ID) code of the System Exclusive data received
by the WSA is different product.''\\[2pt]
\textsc{error 41!} --- ``An error has occured during System Exclusive data reception.The
data from the transmitting device may be incomplete. Please try again.''\\[2pt]
\textsc{error 42!} --- ``An error has occured during System Exclusive transmition.The data
has not been received correctly. Please try again.''
\end{quote}

\estab\ A complete German translation of all three screens exists in \texttt{prom\_b} but
is unreachable: the painter hard-codes the English table, and the language selector byte is
written once, to zero, during reset.
